\originalchapter{半经典近似与三维散射}
\originalsection{WKB的局部构造}
定态一维方程为$-\hbar^2u''/(2m)+Vu=Eu$. 写$u=\exp[\ii S(x)/\hbar]$, 代入得到$(S')^2-\ii\hbar S''=2m(E-V)$. 在允许区令$p(x)=\sqrt{2m(E-V)}>0$, 展开$S=S_0+\hbar S_1+\cdots$. 零阶给$S_0'=\pm p$, 一阶给$2S_0'S_1'=\ii S_0''$, 所以$S_1=(\ii/2)\log p+\text{常数}$. 于是
\[
u_\pm(x)\simeq\frac C{\sqrt{p(x)}}
\exp\left[\pm\frac{\ii}{\hbar}\int^x p(s)\dd s\right].
\]
$1/\sqrt p$不是装饰性的归一化, 它使单向流近似为常数. 近似要求$\hbar|p'|/p^2\ll1$, 即动量在一段局部波长内变化小. 在$E=V$的转折点,$p=0$, 该式发散, 必须另作匹配.

在禁止区令$\kappa(x)=\sqrt{2m(V-E)}/\hbar$, 两个局部解为$\kappa^{-1/2}\exp(\pm\int\kappa\dd x)$. 对一个平滑简单转折点, 将$V-E$线性化后, 方程化为Airy方程. 其衰减解的允许区渐近包含$\pi/4$相位. 两个简单转折点的束缚态匹配给
$\int_{x_1}^{x_2}p(x)\dd x=\pi\hbar(n+1/2)$, $n=0,1,\ldots$.
这里明确使用Airy函数的标准渐近连接公式; 本册完成局部WKB展开, 没有在此证明Airy渐近理论. 硬墙或高阶转折点不能继续使用同一$1/2$修正. 对右转折点$x_2$且右侧禁止, 标准Airy渐近给连接
\[
\frac{C}{\sqrt{|p(x)|}}\exp\!\left[-\frac1\hbar\int_{x_2}^{x}|p(s)|\dd s\right]
\longleftrightarrow
\frac{2C}{\sqrt{p(x)}}\sin\!\left[\frac1\hbar\int_x^{x_2}p(s)\dd s+\frac\pi4\right].
\]
左端同样连接. 两个允许区正弦要表示同一解, 两端$\pi/4$合成$\pi/2$, 给上式$n+1/2$. 此处$|p|=\sqrt{2m(V-E)}$, 与禁止区的$\hbar\kappa$相同.

对谐振子, 转折点$x_\pm=\pm\sqrt{2E/(m\omega^2)}$, 积分是椭圆面积的一半,$\int_{x_-}^{x_+}\sqrt{2mE-m^2\omega^2x^2}\dd x=\pi E/\omega$. 因而WKB给$E=\hbar\omega(n+1/2)$, 恰为精确能级. 这次取等来自二次势的特殊性, 不意味着WKB对任意最低能级都精确.

\originalsection{散射振幅和截面}
短程势的三维弹性散射远处可写为$u(\boldsymbol r)\simeq\e^{\ii kz}+f(\theta,\phi)\e^{\ii kr}/r$, 能量$E=\hbar^2k^2/(2\mu)$. 入射波流密度是$\hbar k/\mu$, 散射球面波的径向流最低阶是$(\hbar k/\mu)|f|^2/r^2$. 在固体角$\dd\Omega$上乘球面面积$r^2\dd\Omega$, 再除入射流密度, 得
$\mathrm d\sigma/\mathrm d\Omega=|f|^2$. 截面量纲为面积, 振幅$f$量纲为长度. 长程库仑势的渐近相位不同, 不能不加修正沿用此短程边界形式.

对中心势, 不同$\ell$通道不混合. $P_\ell$表示勒让德多项式, 径向相移约定为$u_\ell(r)\sim\sin(kr-\ell\pi/2+\delta_\ell)$, 与自由正则波比较. 径向远区相对于自由解多一个相移$\delta_\ell$, 弹性通道幺正性给$S_\ell=\e^{2\ii\delta_\ell}$. 用平面波的球谐展开, 可得
\[
f(\theta)=\frac1k\sum_{\ell=0}^\infty
(2\ell+1)\e^{\ii\delta_\ell}\sin\delta_\ell\,P_\ell(\cos\theta),
\quad
\sigma=\frac{4\pi}{k^2}\sum_\ell(2\ell+1)\sin^2\delta_\ell.
\]
总截面由勒让德多项式正交性积分得到: $\int_{-1}^1P_\ell P_{\ell'}\dd z=2\delta_{\ell\ell'}/(2\ell+1)$. 此处调用平面波的标准球谐展开, 详细相移求解与共振分析仍需对应径向边界问题.

\originalsection{Born近似}
自由外向Green函数为$G(\boldsymbol r-\boldsymbol r')=-\e^{\ii k|\boldsymbol r-\boldsymbol r'|}/(4\pi|\boldsymbol r-\boldsymbol r'|)$, 满足$(\nabla^2+k^2)G=\delta^{(3)}$. 这个负号可由在原点附近积分拉普拉斯检验. 将方程写成$(\nabla^2+k^2)u=(2\mu/\hbar^2)Vu$, 得积分方程
\[
u(\boldsymbol r)=\e^{\ii\boldsymbol k\cdot\boldsymbol r}
-\frac{\mu}{2\pi\hbar^2}\int
\frac{\e^{\ii k|\boldsymbol r-\boldsymbol r'|}}{|\boldsymbol r-\boldsymbol r'|}
V(\boldsymbol r')u(\boldsymbol r')\dd^3r'.
\]
第一Born近似在积分内部用入射波代替完整$u$. 定义末态动量方向$\boldsymbol k'=k\widehat{\boldsymbol r}$. 远区$|\boldsymbol r-\boldsymbol r'|\simeq r-\widehat{\boldsymbol r}\cdot\boldsymbol r'$, 因而
\[
f_{\rm B}(\boldsymbol k',\boldsymbol k)
=-\frac{\mu}{2\pi\hbar^2}\int
\e^{\ii(\boldsymbol k-\boldsymbol k')\cdot\boldsymbol r'}V(\boldsymbol r')\dd^3r',
\quad |\boldsymbol k'|=|\boldsymbol k|.
\]
所以散射测量与势能的Fourier变换相关. 这也说明短程、弱势或合适高能条件为什么有助于近似. 是否足够弱需比较下一次迭代或实际误差; 不能只看到“势很小”就保证在阈值和共振附近有效.


\begin{example}
短程Gaussian势$V(r)=V_*\e^{-r^2/a^2}$, $a>0$. 用Born近似求微分截面.
\end{example}
\begin{solution}
三维Fourier积分可分成三个一维Gaussian, 得$\int\e^{\ii\boldsymbol Q\cdot\boldsymbol r}\e^{-r^2/a^2}\dd^3r=\pi^{3/2}a^3\e^{-a^2Q^2/4}$. 散射角$\theta$对应$Q=|\boldsymbol k-\boldsymbol k'|=2k\sin(\theta/2)$. 所以
$f_B=-\mu\sqrt\pi V_*a^3\e^{-a^2k^2\sin^2(\theta/2)}/(2\hbar^2)$,
$\mathrm d\sigma/\mathrm d\Omega=\mu^2\pi V_*^2a^6\e^{-2a^2k^2\sin^2(\theta/2)}/(4\hbar^4)$.
$\mu V_*a^2/\hbar^2$是本模型的无量纲势强诊断, 但具体误差还依赖$ka$及高阶迭代.
\end{solution}
